% ID: SEL-D0005-OA-E03
% ESTADO: validada. S006, 2026-11-07; tres pasadas del mismo autor y revisión matemática, didáctica y visual.
\textbf{Datos y sucesos.}\par\nopagebreak[4]
Sean $N,P,M$ los sabores naranja, piña y melocotón, y $L$ el suceso ``botella de 2 litros''. Los sabores forman una partición: $P(N)=0.6$, $P(P)=0.3$ y $P(M)=0.1$. En melocotón se da la tasa de 1 litro, de modo que $P(L\mid M)=1-0.6=0.4$.

\textbf{a) Probabilidad total de 2 litros.}\par\nopagebreak[4]
\[
P(L)=0.6(0.8)+0.3(0.7)+0.1(0.4)
=0.48+0.21+0.04=\boxed{0.73}.
\]

\textbf{b) Naranja sabiendo 2 litros.}\par\nopagebreak[4]
Aplicamos Bayes, dividiendo la probabilidad conjunta entre la de la condición:
\[
P(N\mid L)=\frac{P(N)P(L\mid N)}{P(L)}
=\frac{0.48}{0.73}=\boxed{\frac{48}{73}\approx0.657534}.
\]

\textbf{c) Melocotón sabiendo 1 litro.}\par\nopagebreak[4]
Como sólo hay dos capacidades, $L^c$ significa 1 litro y $P(L^c)=0.27$:
\[
P(M\mid L^c)=\frac{0.1(0.6)}{0.27}
=\boxed{\frac29\approx0.222222}.
\]
Comprobación: las probabilidades conjuntas de 1 litro son $0.12,0.09,0.06$; suman 0.27. Junto con $0.48,0.21,0.04$ para 2 litros suman 1. Se eligen botellas al azar, no litros de zumo: no se pondera otra vez por la capacidad.
